\section{Linear Harmonic}
\label{sec:loadcase-linear-harmonic}

\val{LINEARHARMONIC} computes the steady-state response of a linear structural model to a harmonic load over a prescribed frequency sweep. The load amplitudes are defined by the active load collectors and are currently interpreted as real-valued harmonic amplitudes.

\subsection{Governing equation}

At circular excitation frequency \(\omega\), the complex displacement amplitude \(\hat u\) satisfies
\[
\left(K-\omega^2 M+i\omega C\right)\hat u=\hat f .
\]
The complex response is written as
\[
\hat u=u_\mathrm{real}+i u_\mathrm{imag} .
\]
Rayleigh damping uses
\[
C=\alpha M+\beta K .
\]

FEMaster applies the active homogeneous constraints with the Nullspace transformation and solves the reduced frequency-domain problem. To reuse the existing real sparse solver backends, the complex system is written as the real block system
\[
\begin{bmatrix}
A & -B\\
B & A
\end{bmatrix}
\begin{bmatrix}
u_\mathrm{real}\\
u_\mathrm{imag}
\end{bmatrix}
=
\begin{bmatrix}
f\\0
\end{bmatrix},
\]
with
\[
A=K-\omega^2M,
\qquad
B=\omega C .
\]

\subsection{Typical DSL}

\begin{exampledeck}
*LOADCASE, TYPE=LINEARHARMONIC, NAME=FREQUENCY_RESPONSE
  *SUPPORTS
  BC
  *LOADS
  EXCITATION
  *SOLVER, DEVICE=CPU, METHOD=DIRECT
  *CONSTRAINTMETHOD, TYPE=NULLSPACE
  *FREQUENCIES, SCALE=LINEAR
  0.0, 500.0, 201
  *DAMPING, TYPE=RAYLEIGH
  0.0, 1.0e-5
\end{exampledeck}

\cmd{FREQUENCIES} is required and is described in Section~\ref{sec:harmonic-controls}. The current implementation requires the direct solver, Nullspace constraint handling, and homogeneous supports and constraints.

\subsection{Result fields}

For every requested frequency FEMaster writes separate real and imaginary response fields:

\begin{itemize}
\item \val{DISPLACEMENT_REAL} and \val{DISPLACEMENT_IMAG},
\item \val{STRESS_REAL} and \val{STRESS_IMAG},
\item \val{STRAIN_REAL} and \val{STRAIN_IMAG}.
\end{itemize}

The harmonic loadcase passes the excitation frequency to the result writers as the frame value. The native \file{.res} writer currently does not encode this frame value in the field header, so harmonic \file{.res} blocks are associated with frequencies by their sweep order; see the Result Format chapter. Component-wise response magnitude and phase can be recovered from
\[
|\hat u|=\sqrt{u_\mathrm{real}^2+u_\mathrm{imag}^2},
\qquad
\varphi=\operatorname{atan2}(u_\mathrm{imag},u_\mathrm{real}) .
\]
The same interpretation applies to complex stress and strain components.
